\chaves{Engenharia de Requisitos; Sistemas Multiagente; Inteligência Artificial Generativa; Histórias de Usuário; Reconhecimento Automático de Fala}

\begin{resumo}
% TODO: inserir os números quantitativos (WER/CER, F1, alfa de Cronbach,
% cobertura) na frase de resultados, quando os experimentos forem concluídos.
A elicitação de requisitos a partir de reuniões é uma atividade cara, manual e suscetível à perda de informação, na qual o profissional precisa, ao mesmo tempo, participar da discussão e registrá-la em artefatos formais. Este trabalho propõe um sistema multiagente, baseado em inteligência artificial generativa, que automatiza esse processo de ponta a ponta, transformando o áudio de uma reunião em documentação de requisitos estruturada. A arquitetura é composta por quatro agentes especializados, responsáveis, respectivamente, pela transcrição automática do áudio, pela identificação de personas e geração de histórias de usuário segundo o critério INVEST, pela geração do documento de requisitos em dois formatos selecionáveis (o padrão IEEE~830 e um formato específico de cliente) e pela síntese de um diagrama de domínio. Os agentes comunicam-se por meio de um \textit{pipeline} sequencial baseado em arquivos, o que proporciona baixo acoplamento, isolamento de falhas e flexibilidade de ponto de entrada. O sistema foi validado em um caso real, por meio de uma parceria com um órgão público, tendo como referência um documento de requisitos oficial produzido manualmente. A avaliação, de caráter quantitativo, abrange a acurácia da transcrição, a qualidade da extração de informação, a qualidade percebida do documento por profissionais e a propagação de erro ao longo do \textit{pipeline}. % TODO: Os resultados indicam [inserir números].
\end{resumo}
